\chapter{Introducci\'on}
Desde tiempos inmemorables se ha buscado una forma de vincular las ciencias m\'edicas con
el desarrollo tecnol\'ogico, mejorando o creando herramientas digitales que permitan mejores
resultados y ayuden incluso a doctores con gran experiencia. La dermatolog\'ia, rama encargada 
de estudiar enfermedades producidas sobre la superficie de la piel, es una interesante propuesta 
para llevar a cabo esta conexi\'on t\'ecnico-medicinal. El dermat\'ologo debe dar seguimiento
a las lesiones cut\'aneas de sus pacientes, de esa forma si estos poseen una enfermedad 
grave, ser\'a diagnosticada a tiempo, elevando las posibilidades de cura o evoluci\'on positiva.

Para este seguimiento, una herramienta fundamental es el dermatoscopio, un instrumento 
\'optico que permite visualizar estructuras de la piel no visibles a simple vista, 
ampliando las im\'agenes y eliminando el reflejo de la luz superficial. Las im\'agenes 
dermatosc\'opicas resultantes son el insumo principal para el an\'alisis de lesiones 
pigmentadas y no pigmentadas, ya que revelan patrones y caracter\'isticas morfol\'ogicas 
clave (como atipias, redes pigmentarias o vasos anormales) que orientan al especialista 
hacia un diagn\'ostico m\'as certero. De esta forma, la combinaci\'on del dermatoscopio 
con el registro digital de im\'agenes ha abierto la puerta a soluciones automatizadas 
que asisten al m\'edico en la detecci\'on temprana de patolog\'ias como el melanoma.

Detrás de estas soluciones automatizadas se encuentra la inteligencia artificial, 
una disciplina que busca replicar capacidades humanas mediante sistemas computacionales. 
Para comprender el fundamento de estas soluciones y su aplicación en el análisis de imágenes dermatológicas, 
es necesario partir de una definición formal.

Según los autores del libro de referencia en la materia, Stuart Russell y Peter Norvig, 
la Inteligencia Artificial se define como el estudio de agentes computacionales que operan 
de forma autónoma, perciben su entorno, se adaptan al cambio y persiguen objetivos\cite{inteligenciaArtificial}
Como un subcampo de esta, el Machine Learning (o aprendizaje automático) es el estudio 
de algoritmos que permiten a los programas informáticos mejorar automáticamente su 
rendimiento a través de la experiencia\cite{machineLearning}. En esencia, el aprendizaje automático dota a los 
sistemas de la capacidad de aprender sin ser programados explícitamente para cada tarea.

Con el desarrollo y crecimiento de la Inteligencia Artificial (IA), espec\'ificamente el Aprendizaje
de M\'aquinas (ML, por sus siglas en ingl\'es) se han desarrollado diversas herramientas que ayudan 
no solo a dermat\'ologos, sino al personal de salud en general, en sus diagn\'osticos a pacientes.
Hoy en d\'ia el procesamiento de im\'agenes ha abierto nuevas puertas que permiten a los m\'edicos
dar un seguimiento m\'as sencillo y c\'omodo de las distintas lesiones, disminuyendo la fatiga visual y
el uso de equipos externos, adem\'as de ser usado para adiestramiento de futuros doctores. 

\section{Motivaci\'on}
La medicina tiene vital importancia para la sociedad, por ese motivo,
lograr que los m\'edicos tengan una herramienta que les facilite su trabajo,
adem\'as de ser un reto para cualquier persona apasionada por la tecnolog\'ia, constituye una
motivaci\'on tanto social como humana. En el contexto cubano, es de gran ayuda informatizar la mayor
cantidad de procesos, debido a las dificultades log\'isticas que muchas veces se presentan.
Dotar al personal de la salud de una aplicaci\'on web que les permita desde la comodidad de su 
tel\'efono dar seguimiento a sus pacientes y verificar si presentan alguna lesi\'on grave con tan
solo una fotograf\'ia, provocar\'ia un gran avance en el sector medicinal de la naci\'on.
Esa herramienta ya existe desde hace varios a\~nos, concretamente desde 2023 con el desarrollo
del sitio web DermaUH, el cual era capaz de suplir todas estas necesidades mencionadas y brindar las 
comodidades planteadas. Sin embargo, con el alto crecimiento que han tenido las herramientas de IA 
y los modelos de ML, acompa\~nado de la falta de un buen conjunto de im\'agenes de pieles cubanas para entrenar
el modelo desarrollado en su momento, se hace necesario desarrollar una nueva versi\'on de DermaUH y 
vincularle un nuevo modelo m\'as potente con altos porcentajes de efectividad clasificando im\'agenes
de pieles cubanas.

\section{Antecedentes}

El ecosistema DermaUH ha evolucionado a través de dos hitos fundamentales previos a la presente versión 3.0.

\paragraph{Versión 2.0: Sitio web DermaUH (Mauricio Salim Mahmud Sánchez, 2024)}
La versión 2.0 consistió en una plataforma web de apoyo al diagnóstico dermatoscópico, 
desarrollada como una aplicación de una sola página (SPA) con arquitectura cliente-servidor. 
El frontend se construyó con Vue.js y TypeScript, mientras que el backend implementó una API RESTful
construida con Django y Django Rest Framework, utilizando PostgreSQL como sistema gestor de base 
de datos. La plataforma permite a los médicos registrarse (con verificación por administradores), 
gestionar pacientes (CRUD completo), registrar lesiones cutáneas asociadas a cada paciente, y 
almacenar imágenes clínicas y dermatoscópicas por cada lesión. Entre sus funcionalidades avanzadas 
se incluyen: clasificación automática de lesiones mediante un ensemble de redes neuronales 
convolucionales (previamente desarrollado por Claudia Olavarrieta), segmentación de imágenes 
mediante arquitectura U-Net (desarrollada por Marcos Valdivié), y un sistema de colaboración 
que permite compartir lesiones entre especialistas para recibir diagnósticos adicionales. 
La plataforma también expone una API para fines investigativos, accesible por usuarios con 
permisos especiales. Todas las herramientas utilizadas son de código abierto. Esta versión sentó 
las bases del ecosistema, pero dependía completamente de conectividad a internet y estaba 
diseñada exclusivamente para uso en navegador web.

\paragraph{Versión móvil: Aplicación offline con sincronización eventual (Lázaro David Alba Ajete, 2025)}
Posteriormente, se desarrolló una aplicación móvil nativa (Android) que extiende el ecosistema 
al entorno offline. Implementada con Flutter y Dart, sigue los principios de Clean Architecture y 
el patrón MVVM. Utiliza Isar DB como base de datos local embebida y Riverpod para la gestión de estado. 
La aplicación permite a los dermatólogos, en ausencia de conexión a internet, registrar pacientes, 
capturar imágenes de lesiones (clínicas y dermatoscópicas), y ejecutar un modelo de clasificación 
de lesiones (Melanoma, Carcinoma Basocelular, Carcinoma Escamoso) mediante TensorFlow Lite embebido 
en el dispositivo. Los datos se almacenan localmente con campos de sincronización (\texttt{synced}, 
\texttt{deleted}, \texttt{remoteId}) y se transmiten de forma unidireccional al servidor central 
de DermaUH cuando se restablece la conectividad, mediante la API REST preexistente. 
La sincronización puede ser manual o automática diaria configurable. Esta versión resolvió el 
problema de la dependencia de internet, pero su modelo de IA es únicamente clasificador 
(sin capacidades explicativas) y el procesamiento de imágenes se limita a la inferencia del 
modelo previamente entrenado.

\section{Formulaci\'on Del Problema}
DermaUH pretende ser una herramienta potente de apoyo a especialistas en dermatolog\'ia para
el diagn\'ostico de las lesiones m\'as frecuentes de c\'ancer de la piel. La versi\'on 2.0, que
actualmente corre, ha sido evaluada y se han detectado varios problemas que implicar\'an desarrollar
nuevamente el frontend y el backend de la misma. Adem\'as el modelo actual de diagn\'ostico, que 
es un ensemble de 3 redes neuronales no fue entrenado con im\'agenes cubanas.
Por este motivo se desea desarrollar una nueva versi\'on del sitio DermaUH, la cual tendr\'a como 
principal mejora, la incorporaci\'on de nuevos modelos de procesamiento de im\'agenes. A continuaci\'on
se expondr\'an las limitantes y/o deficiencias detectadas en las versiones anteriores del sitio 
web DermaUH.

\begin{enumerate}
    \item El modelo de ML utilizado para clasificar im\'agenes no se entren\'o con im\'agenes de pieles cubanas, alej\'andose en ese aspecto del cliente objetivo, los dermat\'ologos cubanos.
    \item El diagn\'ostico asistido no viene acompa\~nado de una justificaci\'on que responda a la interrogante ¿Por qué se obtuvo este resultado?.
    \item El frontend de la aplicaci\'on tarda alrededor de diez segundos en cargar todos los elementos, haciendo lenta y molesta la experiencia del usuario.
\end{enumerate}

Por lo tanto, el problema que aborda este trabajo se formula de la siguiente manera: 
la versión 2.0 de DermaUH presenta limitaciones en el modelo de clasificación 
(entrenado sin imágenes cubanas), la ausencia de explicabilidad en los diagnósticos, y 
deficiencias de rendimiento en el frontend. A partir de esto, se plantea la siguiente 
pregunta científica:

\begin{quote}
\textbf{¿Es posible desarrollar una versión 3.0 de DermaUH que incorpore un nuevo modelo 
de procesamiento de imágenes (entrenado con pieles cubanas), un modelo explicativo de los 
diagnósticos, y mejore el rendimiento y la seguridad de la plataforma?}
\end{quote}

\section{Objetivos Generales y Espec\'ificos}
Como objetivo general se tendr\'a desarrollar la versi\'on 3.0 de DermaUH, la cual tendr\'a un nuevo
modelo de procesamiento de im\'agenes entrenado con im\'agenes de pieles cubanas, mejorar\'a el rendimiento
de la aplicaci\'on tanto del frontend como el backend, mejorar\'a el protocolo de seguridad y manejo de solicitudes
y notificaciones mediante un cliente de correo Email. Adem\'as esta versi\'on ser\'a validad en 2026
en ensayo cl\'inico. De aqu\'i se derivan varios objetivos espec\'ificos:

\begin{enumerate}
    \item Analizar herramientas existentes para el trabajo con im\'agenes dermatosc\'opicas as\'i como sus ventajas y desventajas para el contexto cubano.
    \item Lograr el desarrollo del frontend y backend que supere las deficiencias presentadas en la v2.0.
    \item Mejorar las m\'etricas de los modelos desarrollados e incorporarlos a la app una vez que sean reentrenados con el nuevo conjunto de im\'agenes de pieles cubanas.
    \item Lograr una app web y explique al especialista la decisi\'on del diagn\'ostico mediante el modelo explicativo (XAI) desarrollado por la Lic.Amanda Noris.
\end{enumerate}

\section{Justificaci\'on}
Se cuenta con una alta capacidad para el desarrollo de la versi\'on 3.0 de DermaUH, debido a todos 
los conocimientos que se tienen sobre el stack tecnol\'ogico seleccionado y los adquiridos durante 
la carrera de Ciencias de la Computaci\'on. Adem\'as de eso existe abundante bibliograf\'ia sobre 
los temas abordados (IA, ML, Procesamiento de im\'agenes), lo cual facilita mucho la b\'usqueda de 
informaci\'on para el desarrollo. Por \'ultimo, la facultad MatCom ya tiene experiencia en el desarrollo
de aplicaciones con un alto componente de IA, sobre todo en aquellas que su principal cualidad
est\'a vinculada al procesamiento de im\'agenes.

\section{Estructura del Trabajo}
El presente trabajo se organiza en cuatro capítulos, además de las conclusiones y recomendaciones finales.

\begin{itemize}
    \item \textbf{Capítulo I. Estado del Arte:} Se analizan sistemas y plataformas existentes para el diagnóstico dermatológico asistido por IA, destacando sus ventajas y limitaciones en el contexto cubano. También se revisan técnicas de procesamiento de imágenes y modelos explicativos (XAI).
    \item \textbf{Capítulo II. Análisis y Diseño:} Se presentan los requisitos funcionales y no funcionales, los diagramas de casos de uso, la arquitectura del sistema (capas, cliente-servidor, API REST) y el diseño detallado de la base de datos.
    \item \textbf{Capítulo III. Implementación:} Se describen las tecnologías utilizadas, la integración de los modelos de IA y procesamiento de imágenes, y las decisiones de implementación más relevantes.
    \item \textbf{Capítulo IV. Pruebas y Resultados:} Se muestran las pruebas funcionales realizadas, las métricas de rendimiento y la validación del sistema con especialistas.
\end{itemize}

Finalmente, se exponen las conclusiones del trabajo y las recomendaciones para futuras líneas de desarrollo.